% Final report for the Coaching Hires Regression Project.
% Compile from this directory twice with: xelatex -interaction=nonstopmode final_report.tex
% final_report.Rmd and final_report.html are retained unchanged from the supplied archive.
\documentclass[11pt,a4paper]{article}
\usepackage[a4paper,margin=1.8cm,headheight=14pt]{geometry}
\usepackage{fontspec}
\setmainfont{Linux Libertine O}
\setsansfont{TeXGyreHeros}
\setmonofont{DejaVu Sans Mono}[Scale=0.82]
\usepackage{amsmath,amssymb,booktabs,tabularx,array,graphicx,float,placeins}
\usepackage{enumitem,microtype,xcolor,fancyhdr,titlesec,caption,hyperref,url,setspace,ragged2e}
\definecolor{navy}{HTML}{17365D}
\definecolor{blue}{HTML}{2F5597}
\hypersetup{colorlinks=true,linkcolor=blue,citecolor=blue,urlcolor=blue,
 pdftitle={Predicting Success of College-Football Coaching Hires},
 pdfauthor={Group RT, BSDS},pdfsubject={Regression analysis and decision framework}}
\graphicspath{{../figures/}}
\setlength{\parindent}{0pt}\setlength{\parskip}{4pt plus 1pt minus 1pt}
\setlength{\tabcolsep}{4pt}\renewcommand{\arraystretch}{1.08}
\setlist[itemize]{leftmargin=1.35em,itemsep=1pt,topsep=2pt}
\setlist[enumerate]{leftmargin=1.6em,itemsep=2pt,topsep=2pt}
\captionsetup{font=small,labelfont=bf,skip=3pt}
\titleformat{\section}{\large\bfseries\color{navy}}{\thesection}{0.5em}{}
\titleformat{\subsection}{\normalsize\bfseries\color{blue}}{\thesubsection}{0.5em}{}
\titlespacing*{\section}{0pt}{9pt}{3pt}\titlespacing*{\subsection}{0pt}{6pt}{2pt}
\pagestyle{fancy}\fancyhf{}
\lhead{\footnotesize Coaching Hires: Regression-Based Decision Framework}
\rhead{\footnotesize Group RT}\cfoot{\footnotesize\thepage}
\renewcommand{\headrulewidth}{0.35pt}\setstretch{1.04}
\newcolumntype{Y}{>{\RaggedRight\arraybackslash}X}
\newcommand{\cpi}{\textsc{CPI}}
\newcommand{\tabnote}[1]{\vspace{2pt}{\parbox{\linewidth}{\footnotesize\textit{Note.} #1}}}

\begin{document}

% TITLE PAGE: included in the 20-page limit
\thispagestyle{empty}
\vspace*{0.65cm}
{\color{navy}\Large\bfseries BSDS PROJECT REPORT}\par
\vspace{0.65cm}
{\color{navy}\Huge\bfseries Predicting Success of College-Football Coaching Hires}\par
\vspace{0.25cm}
{\Large A Regression-Based Decision Framework}\par
\vspace{0.55cm}\hrule\vspace{0.45cm}
\begin{tabularx}{\linewidth}{@{}>{\bfseries}p{3.2cm}Y@{}}
Programme & BSDS\\
Group folder & RT (replace with the official group name if different)\\
Research question & Do observable pre-hire characteristics predict subsequent team performance?\\
Primary outcome & Project-defined Coach Performance Index: team SP+ relative to FBS mean (\verb|sp_rel_team|)\\
Data window & Hires beginning 1999--2026; completed-hire outcomes observed through 2025\\
Report date & 3 November 2026\\
\end{tabularx}
\vspace{0.45cm}\hrule\vspace{0.35cm}
\section*{Abstract}
This project examines whether observable pre-hire characteristics predict the subsequent on-field performance of college-football head coaches. The supplied team-compiled dataset contains 132 hire records and 60 fields. The primary regression analysis uses 84 eligible completed hires; the survival analysis uses 120 completed hires; and a separate temporal holdout evaluates 42 hires beginning in 2022. We distinguish a 15-slope confirmatory linear model, retained to test four stated hypotheses, from a parsimonious Bayesian information criterion (BIC)-selected model used for prediction. The confirmatory model explains 24.7\% of in-sample variation but has adjusted \(R^2=0.081\) and 10-fold cross-validated \(R^2=-0.159\), worse than the null benchmark. The BIC-selected two-predictor model reaches cross-validated \(R^2=0.071\), a modest result. Prior offensive-coordinator experience is nominally associated with higher CPI, but none of the four primary hypotheses survives Holm correction. The 2022+ holdout does not demonstrate improvement over a training-mean baseline. Mixed-effects and holdout result discrepancies are disclosed rather than silently resolved.

\textbf{Overall conclusion.} The analysis provides a transparent description of candidate signals in this dataset, but it does not establish a validated predictor of coaching success. The Coach Hire Score is exploratory and should not be used as a stand-alone hiring rule.
\vfill
\small The report is within the professor's 20-page limit including references and appendices. The companion archives contain the raw dataset, processed data, R scripts, helper code, figures, full numerical result tables, and supporting records.
\newpage

\tableofcontents
\newpage

\section{Introduction and Research Questions}
College-football head-coaching appointments are high-stakes decisions made under uncertainty. A candidate's prior record, coordinator experience, playing background, relationships with the hiring institution, and the quality of the inherited programme may all be relevant. The statistical task is not merely to rank biographies retrospectively, but to assess whether pre-hire information has a stable association with a clearly defined post-hire performance outcome.

The primary research question is: \textbf{How well do observable pre-hire characteristics predict the project-defined Coach Performance Index (CPI) of a college-football head-coaching hire?} The response is \verb|sp_rel_team|, defined in the project data dictionary as the team's SP+ rating relative to the FBS mean. A win-percentage differential is retained as a secondary outcome for a holdout sensitivity analysis.

\begin{table}[H]\centering\small
\caption{Primary hypotheses specified for confirmatory inference}
\begin{tabularx}{\linewidth}{@{}lYlY@{}}
\toprule Hypothesis & Predictor & Expected sign & Interpretation\\\midrule
H1 & Prior head-coaching win percentage (\verb|AllPriorWinPct|) & Positive & Stronger prior HC record predicts higher CPI.\\
H2 & Connection to hiring school (\verb|ConnectionToHiring.School|) & Positive & Prior institutional connection predicts higher CPI.\\
H3 & Prior offensive-coordinator experience (\verb|PreviousOC|) & Positive & Previous OC experience predicts higher CPI.\\
H4 & NFL playing experience (\verb|PlayedNFL|) & Positive & Prior NFL playing experience predicts higher CPI.\\
\bottomrule\end{tabularx}\end{table}

Additional analyses compare model-selection methods, assess survival and retention outcomes, and develop an interpretable candidate-ranking score. These secondary exercises do not replace the confirmatory H1--H4 tests.

\section{Data, Sources, and Sample Construction}
\subsection{Dataset and provenance}
The raw dataset is a project-compiled CSV with \textbf{132 rows and 60 columns}. It contains coach identity and career-history variables, hiring-school and season identifiers, programme-performance measures, and outcome indicators. The \verb|Wikipedia| column stores a coach-profile URL for every record (110 unique URLs across the 132 rows). The team assembled the season-level SP+ and win-percentage quantities used to construct the outcome and programme-context variables.

\textbf{Downloadable data source.} The exact CSV used by the analysis is distributed alongside this report at \href{../data/raw/data.csv}{\texttt{analysis/data/raw/data.csv}} and is included in both ZIP archives. Its SHA-256 digest is \texttt{059dc4c7186128d431ade38213f70e3bf044f3a18f99e466704cb12fe63a3d81}, recorded in \verb|data/raw/data.csv.sha256|. This is a team-compiled dataset; the report does not claim that one public downloadable source file exists. Coach-profile links are preserved at record level, but a complete source ledger mapping every historical SP+ input to its original publication is absent. This is a material data-provenance limitation and should be corrected before external replication or consequential use.

\subsection{Outcome, predictors, and sample flow}
The primary response \verb|sp_rel_team| is the programme's SP+ relative to the FBS mean. \verb|team_sp_recent| is prior programme SP+ context, while \verb|sp_prev_rel| describes the coach's prior-role SP+ relative to the FBS mean. Other predictors include prior HC status, prior win percentages, coordinator roles, playing/NFL coaching experience, age at hire, conference movement, prior championship, and internal promotion. \verb|winpct_rel_team| is the secondary outcome. Definitions are supplied in \verb|data/README.md|.

\begin{table}[H]\centering\small
\caption{Analysis populations and their distinct purposes}
\begin{tabularx}{\linewidth}{@{}l r Y@{}}\toprule
Population & \(n\) & Construction and use\\\midrule
Raw file & 132 & All compiled hire records, including 12 speculative 2026 hires.\\
Completed-hire frame & 120 & A last season is observed; used for tenure/survival analysis.\\
Regression sample & 84 & At least two completed seasons, project employment-status filter, and Sam Pittman excluded for a prior-HC coding inconsistency.\\
Temporal holdout & 42 & Hires beginning in 2022 or later with observed CPI.\\
Holdout training set & 75 & Regression records after removing holdout identities.\\\bottomrule
\end{tabularx}\end{table}

The original plan specified at least three completed seasons; this was relaxed to two to preserve the planned minimum analysis size. The regression construction also retains a hire when \verb|StillEmployedAtSchool = 0| or \verb|first_season| is at most 2021. This condition can select recent hires based on observed tenure/employment status and is therefore a source of selection bias, not a neutral restriction. The 12 speculative 2026 hires have no observed last season and are excluded from analyses requiring completed outcomes.

\subsection{Missing data and integrity checks}
The raw file has 34 missing values in \verb|AllPriorWinPct|, largely structural for first-time head coaches; these are set to 0.50 as a neutral placeholder. Two anomalous prior-HC cases (Chad Morris and Dana Holgorsen) are filled from \verb|PreviousWinPct|. Sam Pittman is excluded because the file flags him as a prior head coach while recording zero prior wins and losses and a 0.50 no-prior-HC placeholder; imputing from that value would be circular. Missing \verb|sp_prev_rel| values (37 in the raw file) are imputed to the regression-sample mean. These deterministic imputations are pragmatic and do not propagate imputation uncertainty.

Five records are flagged in \verb|tables/data_corrections.csv|. Hugh Freeze and Luke Fickell have internal contradictions between firing and employment-status indicators. Brian Kelly, Mike Gundy, and James Franklin have externally flagged employment-status inconsistencies. The primary data retain these records as coded; a survival sensitivity analysis excludes the three externally flagged coaches. Full row-level flags and inclusion/exclusion lists are supplied.

\section{Statistical Methods}
\subsection{Confirmatory linear model}
The inferential specification is an ordinary least-squares model with 15 slope terms and an intercept (16 fitted parameters):
\begin{align*}
\mathrm{CPI}_i={}&\beta_0+\beta_1\mathrm{PrevHC}_i+\beta_2\mathrm{AllPriorWinPct}_i+\beta_3\mathrm{PreviousWinPct}_i\\
&+\beta_4\mathrm{PreviousOC}_i+\beta_5\mathrm{PreviousDC}_i+\beta_6\mathrm{SideOfBall}_i+\beta_7\mathrm{Connection}_i\\
&+\beta_8\mathrm{PlayedNFL}_i+\beta_9\mathrm{Age50}_i+\beta_{10}\mathrm{NFLCoaching}_i\\
&+\beta_{11}(\mathrm{PrevHC}_i\times\mathrm{PriorSP+}_i)+\beta_{12}\mathrm{WonCC}_i+\beta_{13}\mathrm{LeavingP4}_i\\
&+\beta_{14}\mathrm{RecentTeamSP+}_i+\beta_{15}\mathrm{PromoteWithin}_i+\epsilon_i.
\end{align*}
Binary predictors use 0 as the reference level. This full model is retained for the stated H1--H4 inference even though it is not the strongest predictive model. H1 is operationalised as the net contrast \(H_0:\beta_{\mathrm{AllPriorWinPct}}+\beta_{\mathrm{PreviousWinPct}}=0\), because the two win-percentage variables overlap conceptually and have opposite fitted signs.

\subsection{Model selection and validation}
A bidirectional BIC search starts from the null model and considers the full specification as its upper scope, using penalty \(k=\log(n)\). AIC stepwise, best-subset criteria, ridge, principal-components regression (PCR), LASSO, and elastic net provide secondary checks. Ten-fold cross-validation uses fixed seed 2026 and common folds for the principal model comparison. The BIC-selected model is used for prediction; the full model is used for confirmatory inference. This distinction avoids treating a post-selection specification as if it had been pre-specified for hypothesis testing [1,2].

\subsection{Diagnostics and secondary methods}
Residual-versus-fitted, Q--Q, scale-location, Cook's distance, leverage, and DFBETA diagnostics were examined. Breusch--Pagan and White squares-only tests assess heteroscedasticity; a Box--Cox analysis assesses response transformation. WLS and regularised models are sensitivity analyses. Survival duration is last season minus first season plus one; on-field firing is the event, while off-field firing is censored. Kaplan--Meier curves and a Cox model assess tenure [3,4]. A logistic model examines a four-season retention indicator, but its outcome coding does not distinguish insufficient follow-up and is therefore secondary only. Bootstrap and jackknife calculations assess uncertainty in the prior-win coefficient [5].

\section{Descriptive Results and Regression Diagnostics}
The regression sample has mean CPI approximately -1.96 and standard deviation 8.44; mean prior career HC win percentage is 0.637 (SD 0.167), mean previous-job win percentage is 0.579 (SD 0.119), and mean prior programme SP+ is 7.83 (SD 8.91). These are descriptive summaries. Programme resources, schedule, conference context, institutional investment, and time can affect both hire characteristics and outcomes.

\begin{figure}[H]\centering\includegraphics[width=0.76\linewidth]{fig1_cpi_by_hire_type.png}
\caption{Observed CPI distribution by hire type. Differences are descriptive, not adjusted causal effects.}\end{figure}

\subsection{Residual and variance diagnostics}
\begin{figure}[H]\centering\includegraphics[width=0.82\linewidth]{fig3_residual_panel.png}
\caption{Confirmatory OLS diagnostics: residuals versus fitted values, normal Q--Q plot, scale-location, and Cook's distance.}\end{figure}

The archived tests show no strong evidence of heteroscedasticity: Breusch--Pagan \(\chi^2=15.76\), 15 df, \(p=0.398\); White squares-only \(\chi^2=22.56\), 30 df, \(p=0.833\). Box--Cox gives \(\hat\lambda=0.909\), with a 95\% interval including 1, so the original response scale is retained. Durbin--Watson is 2.132 (\(p=0.736\)); because observations are cross-sectional hires rather than a naturally ordered time series, this test is descriptive and does not establish independence.

Influence flags are not negligible: 8 observations exceed Cook's-distance rule \(4/n\), 3 exceed leverage rule \(2p/n\), and 33 exceed the stored DFBETA threshold. No observation was deleted solely on these flags. Maximum Cook's distance is 0.112. The largest VIF is 5.55 for \verb|PreviousDC|; VIFs for \verb|AllPriorWinPct| and \verb|PreviousWinPct| are 4.11 and 3.59. No VIF exceeds 10, but the overlapping prior-win variables and their opposite coefficient signs make the net H1 contrast more interpretable than either partial coefficient alone.

\section{Model Comparison and Primary Findings}
\subsection{Predictive model comparison}
\begin{table}[H]\centering\small
\caption{Model comparison on the 84-observation regression sample}
\begin{tabular}{@{}lrrrrrr@{}}\toprule
Model & Parameters & \(R^2\) & Adj. \(R^2\) & AIC & BIC & CV \(R^2\)\\\midrule
Null & 1 & 0.000 & 0.000 & 599.62 & 604.49 & -0.056\\
PreviousOC only & 2 & 0.100 & 0.089 & 592.73 & 600.03 & 0.047\\
BIC-selected & 3 & 0.154 & 0.133 & 589.55 & 599.27 & 0.071\\
Confirmatory full & 16 & 0.247 & 0.081 & 605.77 & 647.09 & -0.159\\\bottomrule\end{tabular}
\tabnote{Ten-fold folds use seed 2026. Cross-validated performance is the relevant predictive comparison; high in-sample fit does not establish generalisation.}\end{table}

The BIC-selected specification is \(\widehat{\mathrm{CPI}}=-1.438+8.516\,\mathrm{PreviousOC}-0.223\,\mathrm{team\_sp\_recent}\). The PreviousOC coefficient is positive (SE 2.482; 95\% CI [3.578, 13.454]; nominal \(p=0.001\)); recent-team SP+ is negative (SE 0.098; 95\% CI [-0.418, -0.028]; nominal \(p=0.026\)). These are post-selection estimates and should be considered exploratory. The negative context coefficient is compatible with regression-to-the-mean patterns but is not a causal effect.

\begin{figure}[H]\centering\includegraphics[width=0.72\linewidth]{fig5_model_selection_profile.png}
\caption{Model-selection profile. BIC favours a compact specification over the 16-parameter confirmatory model.}\end{figure}

Ridge and PCR improve on the full OLS cross-validation RMSE in the archived sensitivity run: ridge (lambda-min) RMSE 8.329 and PCR with nine components RMSE 8.416, versus full OLS RMSE 9.027. LASSO/elastic-net outputs also concentrate selection on PreviousOC and recent programme SP+ at selected penalty levels, but selected terms vary with tuning. Four-block F tests do not show any predictor block surviving Holm adjustment; the coordinator-background block has raw \(p=0.106\), adjusted \(p=0.422\). These checks favour parsimony but do not establish a stable predictive model.

\subsection{Confirmatory tests of H1--H4}
\begin{table}[H]\centering\small
\caption{Confirmatory H1--H4 coefficient tests from the full model}
\begin{tabular}{@{}llrrrrl@{}}\toprule
Hypothesis & Term & Estimate & SE & Raw \(p\) & Holm \(p\) & Verdict\\\midrule
H1 & AllPriorWinPct & 7.720 & 12.332 & 0.533 & 1.000 & Not supported\\
H2 & ConnectionToHiring.School & -0.425 & 2.409 & 0.860 & 1.000 & Not supported\\
H3 & PreviousOC & 7.966 & 3.853 & 0.043 & 0.170 & Not after correction\\
H4 & PlayedNFL & -3.947 & 3.197 & 0.221 & 0.664 & Not supported\\\bottomrule\end{tabular}
\tabnote{Holm step-down correction across four stated tests. Bonferroni-adjusted p-values are included in the archived table. No hypothesis survives multiplicity adjustment at 5\% [7].}\end{table}

H3 is the only nominally significant coefficient before correction. Its positive estimate is compatible with the hypothesised direction, but its Holm-adjusted p-value is 0.170 and it is not confirmatory evidence. H1 is parameterisation-sensitive: the two prior-win coefficients are +7.720 and -13.489, giving net estimate -5.769. The net-effect test gives \(F(1,68)=0.215\), \(p=0.644\), and 95\% CI [-30.578, 19.040]. A restricted model dropping \verb|PreviousWinPct| gives an \verb|AllPriorWinPct| estimate of 1.060 (SE 10.175; \(p=0.917\)). Neither supports H1. The full-model overall F test is not significant (\(F=1.489\), 15 and 68 df, \(p=0.134\)).

\begin{table}[H]\centering\small
\caption{BIC-selected model coefficients}
\begin{tabular}{@{}lrrrr@{}}\toprule
Term & Estimate & SE & 95\% CI & Raw \(p\)\\\midrule
Intercept & -1.438 & 1.166 & [-3.757, 0.882] & 0.221\\
PreviousOC & 8.516 & 2.482 & [3.578, 13.454] & 0.001\\
team\_sp\_recent & -0.223 & 0.098 & [-0.418, -0.028] & 0.026\\\bottomrule\end{tabular}
\tabnote{Selected on the full 84-row regression sample. In holdout evaluation, the BIC selection is repeated using the 75-row training set only.}\end{table}

\section{Survival, Mixed-Effects, and Retention Analyses}
\subsection{Tenure and firing}
The survival frame contains 120 completed hires and 39 on-field firing events. Tenure is calculated as last season minus first season plus one; off-field firings are censored rather than treated as on-field-performance events. Kaplan--Meier curves by prior HC experience are similar, and the log-rank test is not significant (\(\chi^2=1.571\), \(p=0.210\)).

\begin{figure}[H]\centering\includegraphics[width=0.69\linewidth]{fig2_km_curves.png}
\caption{Kaplan--Meier tenure curves by prior head-coaching experience.}\end{figure}

The Cox model estimates HR 1.752 for prior win percentage (95\% CI [0.189, 16.259], \(p=0.622\)), HR 1.217 for prior HC experience (95\% CI [0.535, 2.768], \(p=0.640\)), and HR 1.032 per year for age at hire (95\% CI [0.981, 1.086], \(p=0.223\)). All intervals include 1; the global proportional-hazards test has \(p=0.139\). Excluding the three externally flagged coaches leaves 117 hires and 37 events; log-rank \(p=0.226\), with prior-win and prior-HC hazard ratios 1.718 and 1.183. With only 39 events, these results are underpowered and inconclusive rather than evidence of no association.

\subsection{School and conference random effects}
The archived mixed-effects outputs are not fully reconciled. \verb|w5_mixed_summary.csv| reports school ICC 0.041 and LRT \(p=0.893\), while the later ML-only correction in \verb|lrt_school_random.csv| reports a boundary statistic of zero, \(p=0.500\), and ICC 0.000. The earlier \verb|09_mixed_effects.R| compares a GLS model fit by ML with an LME model fit by REML, so the likelihoods are not a clean basis for an LRT comparison; the later \verb|19_lrt_fix.R| uses ML for both fits with a boundary correction [6]. Because the saved ICC summaries still disagree, this report does not claim school-level heterogeneity is exactly zero. The defensible conclusion is that the archived outputs provide no reliable evidence of a material school random intercept; the fixed-effects model remains primary pending a clean rerun and reconciliation.

\subsection{Logistic retention analysis}
The secondary logistic model defines \verb|Retained4| as at least four completed seasons without an off-field firing. Odds-ratio intervals are very wide (AllPriorWinPct OR 9.659, 95\% CI [0.127, 732.056]; PreviousOC OR 0.770, 95\% CI [0.179, 3.304]); no displayed predictor interval excludes 1. The Hosmer--Lemeshow test has \(p=0.433\), but a non-significant calibration test does not establish discrimination. More importantly, the coding classifies not-yet-observable four-season cases as non-retained, conflating censoring with failure. These results are descriptive only and do not modify H1--H4.

\section{Resampling, Temporal Validation, and Decision Score}
\subsection{Bootstrap and jackknife}
The 1,000-resample bootstrap interval for the AllPriorWinPct coefficient is [-19.86, 34.22], including zero. Jackknife bias is 0.34, small relative to the reported standard error of approximately 13.3. The resampling checks confirm that the prior-win association is imprecise; they do not resolve confounding or collinearity.

\subsection{2022+ holdout}
The holdout script defines the test set as all observed-CPI hires beginning in 2022 or later, without filtering on current employment status. Training is the regression sample after removing holdout identities (75 observations). Both models are refitted on training data only. The BIC model gives 95\% prediction-interval coverage 0.952 (40/42; exact 95\% CI [0.838, 0.994]), RMSE 9.004 and MAE 6.934. The train-only confirmatory model gives the same coverage, RMSE 9.159 and MAE 7.103.

\begin{figure}[H]\centering\includegraphics[width=0.72\linewidth]{fig7_pi_coverage.png}
\caption{Observed CPI and 95\% prediction intervals for the 2022+ holdout under the train-only BIC model.}\end{figure}

This is \textbf{post hoc temporal validation}, not a pre-registered prospective test. The original nine-row subset used a \verb|StillEmployedAtSchool = 0| condition, consisted only of fired/departed coaches, and overlapped the training sample. Its coverage (0.889; exact interval [0.518, 0.997]) is retained in the archive for traceability but is not the primary validation. The 42-row cohort corrects the status filter and removes overlap, but the decision to replace the initial subset was made after its result was observed; this limits confirmatory claims.

The separate \verb|decision_rule_validation.csv| reports Spearman \(\rho=0.189\), RMSE 8.952, and training-mean baseline RMSE 8.874; it records zero cases achieving a 10\% improvement over the baseline. The RMSE of 8.952 differs slightly from 9.004 in \verb|holdout_performance.csv|. This discrepancy is unresolved in the supplied outputs. Both show that prediction is weak and does not beat the simple baseline. For the win-percentage-differential outcome, holdout coverage is 0.881 and RMSE 0.192; this RMSE is on a different scale and is not comparable to CPI RMSE.

\subsection{Coach Hire Score}
The score maps a model fitted value to its percentile within the training fitted-value distribution; it is not a probability of success. In the worked examples, BIC-model percentile scores are 45.2 for an established head coach, 96.4 for an experienced OC, and 97.6 for a first-time HC profile. Their BIC prediction intervals are approximately 31--33 CPI points wide and overlap substantially. The ranking is heavily influenced by prior OC status and programme context. A percentile of 96.4 does not mean a 96.4\% chance of success. Given the holdout rank correlation of 0.189 and RMSE worse than the mean baseline, the score is an exploratory discussion aid, not an operational screening or hiring rule.

\section{Discussion and Limitations}
The main finding is the gap between in-sample fit and predictive performance. The confirmatory model reaches \(R^2=0.247\), but adjusted \(R^2=0.081\) and cross-validated \(R^2=-0.159\) indicate poor generalisation. The BIC model is more economical and performs better under cross-validation, but CV \(R^2=0.071\) is modest and the post hoc holdout does not demonstrate improvement over a training-mean benchmark. H3 is a plausible exploratory association, not a confirmed result after multiplicity adjustment.

Material limitations are:
\begin{enumerate}
\item \textbf{Small, selected sample:} 84 regression observations for 16 fitted parameters, with two-season and employment-status restrictions. The employment-status condition may select recent hires on observed tenure/outcome.
\item \textbf{Observational confounding:} resources, schedules, institutional investment, conference context, interim staff, programme trajectory, and timing may affect both hire attributes and outcomes. Coefficients are associations, not causal effects.
\item \textbf{Collinearity and multiplicity:} the two prior-win measures overlap; H1 depends on a net contrast. No hypothesis survives Holm adjustment.
\item \textbf{Source provenance:} coach-profile links are embedded in the CSV, but a complete citation trail for historical SP+ input values is missing. The CSV is distributed, but historical input sourcing is not fully auditable.
\item \textbf{Post hoc holdout:} the 42-row temporal cohort corrects a status-filter/overlap problem but was chosen after observing the initial nine-row result. It is not an untouched external test set.
\item \textbf{Survival/retention limitations:} only 39 on-field firing events support the Cox model; the Retained4 coding does not distinguish insufficient follow-up from non-retention.
\item \textbf{Archived output conflicts:} mixed-effects ICC/LRT files disagree and the two holdout RMSE outputs differ slightly. They are disclosed, not averaged or silently resolved.
\item \textbf{No operational threshold:} neither causal effects nor sufficiently precise prediction intervals are established for candidate selection.
\end{enumerate}

\section{Conclusion}
In this team-compiled sample, prior offensive-coordinator experience is the strongest candidate signal and recent programme SP+ also enters the BIC-selected model. However, no pre-specified hypothesis survives Holm correction, the full model performs worse than the null under cross-validation, and the selected model's temporal holdout errors do not beat a simple mean baseline. Survival findings are imprecise, and the retention model has an outcome-censoring limitation. The project supports a transparent statistical description of the available data, not a validated predictor of coaching success. Further work should prioritise a full source ledger, an untouched pre-specified temporal test set, more hires and seasons, censoring-aware retention outcomes, and reconciliation of mixed-model results before the score is considered for practical use.

\section{Reproducibility and Submission Contents}
The project includes 16 main R scripts in \verb|analysis/R/| and helper scripts in \verb|analysis/lib/|. The pipeline runner is \verb|analysis/run_all.sh| and sets seed 2026. The submission archive contains the raw CSV and checksum, six processed RDS files, all main R scripts and helper functions, 19 figures, 55 result tables, weekly memos, audit/determinism documentation, session information, and this final report. The supplied \verb|final_report.Rmd| and \verb|final_report.html| are preserved byte-for-byte unchanged.

\textbf{Build instructions.} From \verb|RT/analysis/reports/| run \verb|xelatex -interaction=nonstopmode final_report.tex| twice. The PDF links to the local source dataset at \verb|../data/raw/data.csv|. Detailed results and inclusion/exclusion records are supplied as companion CSV files so the report can stay within the 20-page limit.

\textbf{Environment and verification.} The project records R 4.6.1 on Linux x86\_64 and timezone Asia/Kolkata. The preparation container did not have \verb|Rscript|, so the archived R pipeline could not be freshly rerun here. This PDF was compiled with XeLaTeX and checked for page count and layout; numerical findings were taken from the supplied result CSVs and existing R Markdown output. Before claiming a fresh end-to-end rebuild, run \verb|run_all.sh| in the target R environment and reconcile the two explicitly identified output discrepancies.


\appendix
\section{Complete Confirmatory OLS Coefficient Table}
The following table preserves the full coefficient output for the 15-slope confirmatory specification, including confidence intervals. The table is included to make the confirmatory model auditable without reproducing every raw R output in the main text.
\begin{table}[H]\centering\scriptsize
\caption{Full confirmatory OLS model: coefficient estimates, standard errors, p-values, and 95\% confidence intervals}
\begin{tabularx}{\linewidth}{@{}lrrrrY@{}}\toprule
Term & Estimate & SE & \(p\) & 95\% CI\\\midrule
Intercept & 0.579 & 8.450 & 0.946 & [-16.282, 17.440]\\
PrevHCAnyLevel & -1.050 & 3.930 & 0.790 & [-8.893, 6.793]\\
AllPriorWinPct & 7.720 & 12.332 & 0.533 & [-16.888, 32.328]\\
PreviousWinPct & -13.489 & 14.093 & 0.342 & [-41.611, 14.633]\\
PreviousOC & 7.966 & 3.853 & 0.043 & [0.276, 15.655]\\
PreviousDC & 0.189 & 4.599 & 0.967 & [-8.989, 9.366]\\
SideOfBallO & -0.513 & 3.149 & 0.871 & [-6.797, 5.771]\\
ConnectionToHiring.School & -0.425 & 2.409 & 0.860 & [-5.232, 4.382]\\
PlayedNFL & -3.947 & 3.197 & 0.221 & [-10.325, 2.432]\\
AgeAtHiringMinus50 & -0.150 & 0.170 & 0.380 & [-0.489, 0.189]\\
NFLCoachingExperience & 0.796 & 2.244 & 0.724 & [-3.682, 5.274]\\
WonCCasHCany & 4.514 & 2.960 & 0.132 & [-1.393, 10.421]\\
Conf.LeavingP4 & -0.300 & 2.488 & 0.904 & [-5.265, 4.664]\\
team\_sp\_recent & -0.266 & 0.117 & 0.026 & [-0.499, -0.033]\\
PromoteWithin & 3.021 & 3.673 & 0.414 & [-4.308, 10.350]\\
PrevHCAnyLevel:sp\_prev\_rel & 0.083 & 0.170 & 0.629 & [-0.257, 0.423]\\
\bottomrule\end{tabularx}
\tabnote{The nominal PreviousOC and team\_sp\_recent p-values do not account for the full model-selection process. H1--H4 multiplicity-adjusted tests are shown in the main report.}
\end{table}

\section{Supplementary Survival and Logistic Outputs}
\begin{table}[H]\centering\small
\caption{Cox proportional-hazards estimates}
\begin{tabular}{@{}lrrrr@{}}\toprule
Term & Hazard ratio & 95\% CI lower & 95\% CI upper & \(p\)\\\midrule
AllPriorWinPct & 1.752 & 0.189 & 16.259 & 0.622\\
PrevHCAnyLevel & 1.217 & 0.535 & 2.768 & 0.640\\
AgeAtHiringMinus50 & 1.032 & 0.981 & 1.086 & 0.223\\\bottomrule\end{tabular}
\end{table}
\begin{table}[H]\centering\small
\caption{Secondary logistic-regression odds ratios for Retained4}
\begin{tabular}{@{}lrrr@{}}\toprule
Term & OR & 95\% CI lower & 95\% CI upper\\\midrule
AllPriorWinPct & 9.659 & 0.127 & 732.056\\
PreviousOC & 0.770 & 0.179 & 3.304\\
PreviousDC & 1.434 & 0.436 & 4.711\\
ConnectionToHiring.School & 2.095 & 0.702 & 6.254\\
PlayedNFL & 0.875 & 0.190 & 4.025\\
sp\_prev\_rel & 0.995 & 0.921 & 1.076\\\bottomrule\end{tabular}
\tabnote{Intervals are very wide and all include 1. The Retained4 coding limitation is discussed in Section 6.}
\end{table}

\section{Supporting-Artifact Index}
\begin{table}[H]\centering\small
\caption{Where to find the complete supporting materials in the submission folder}
\begin{tabularx}{\linewidth}{@{}p{3.4cm}Y@{}}\toprule
Artifact & Relative path within the group folder\\\midrule
Raw downloadable dataset & \verb|analysis/data/raw/data.csv| and checksum \verb|data.csv.sha256|\\
Processed datasets/models & \verb|analysis/data/processed/*.rds|\\
Analysis code & All 16 scripts in \verb|analysis/R/|\\
Shared model/pipeline helpers & \verb|analysis/lib/|\\
Reproducible runner & \verb|analysis/run_all.sh|\\
Report source and outputs & \verb|analysis/reports/final_report.tex| and \verb|final_report.pdf|\\
Preserved report render & \verb|analysis/reports/final_report.Rmd| and \verb|final_report.html| (unchanged)\\
Figures and full numerical outputs & \verb|analysis/figures/| and \verb|analysis/tables/|\\
Memos and run provenance & \verb|analysis/reports/memos/|, \verb|session_info.txt|, \verb|determinism_report.md|, \verb|audit_log.csv|\\
\bottomrule\end{tabularx}
\end{table}

\section*{References}
\begin{enumerate}[label={[\arabic*]},leftmargin=2.1em,itemsep=2pt]
\item Burnham, K. P., and Anderson, D. R. (2002). \textit{Model Selection and Multimodel Inference: A Practical Information-Theoretic Approach}. 2nd ed. Springer.
\item James, G., Witten, D., Hastie, T., and Tibshirani, R. (2021). \textit{An Introduction to Statistical Learning: With Applications in R}. 2nd ed. Springer.
\item Kaplan, E. L., and Meier, P. (1958). Nonparametric estimation from incomplete observations. \textit{Journal of the American Statistical Association}, 53(282), 457--481.
\item Cox, D. R. (1972). Regression models and life-tables. \textit{Journal of the Royal Statistical Society: Series B}, 34(2), 187--220.
\item Efron, B., and Tibshirani, R. J. (1993). \textit{An Introduction to the Bootstrap}. Chapman \& Hall/CRC.
\item Self, S. G., and Liang, K.-Y. (1987). Asymptotic properties of maximum likelihood estimators and likelihood ratio tests under nonstandard conditions. \textit{Journal of the American Statistical Association}, 82(398), 605--610.
\item Holm, S. (1979). A simple sequentially rejective multiple test procedure. \textit{Scandinavian Journal of Statistics}, 6(2), 65--70.
\item White, H. (1980). A heteroskedasticity-consistent covariance matrix estimator and a direct test for heteroskedasticity. \textit{Econometrica}, 48(4), 817--838.
\item Breusch, T. S., and Pagan, A. R. (1979). A simple test for heteroscedasticity and random coefficient variation. \textit{Econometrica}, 47(5), 1287--1294.
\item Group RT (2026). \textit{Coaching Hires Dataset}, 132 records x 60 fields, team-compiled CSV. Included in the companion submission files; SHA-256 recorded in \verb|data.csv.sha256|.
\end{enumerate}

\end{document}
